\section*{Chapter \ref{ch:b1:isostasy-gravity} --- Isostasy and Gravity Anomalies}

\begin{solution}{exo:b1:isostasy-gravity:1}
(a) $\qty{-150}{\milli\gal} = -150\times\qty{e-5}{m.s^{-2}} = \qty{-1.5e-3}{m.s^{-2}}$,
and $\num{1.5e-3}/9.81 = \num{1.53e-4}$ of $g$. (b) $\qty{1}{\milli\gal} =
\qty{e-5}{m.s^{-2}}$, \num{1.02e-6} of $g$; $\qty{1}{\micro\gal} =
\qty{e-8}{m.s^{-2}}$, \num{1.02e-9} of $g$: a gravimeter that resolves a
microgal measures gravity to about one part in a billion.
\end{solution}

\begin{solution}{exo:b1:isostasy-gravity:2}
(a) At the equator both sines vanish: $\gamma = \qty{9.780327}{m.s^{-2}}$. At
\ang{45}, $\sin^2\varphi = 0.5$ and $\sin^2 2\varphi = 1$:
$\gamma = 9.780327\times(1 + 0.0026512 - 0.0000058) = \qty{9.806200}{m.s^{-2}}$.
At the poles, $\sin^2\varphi = 1$ and $\sin 2\varphi = 0$:
$\gamma = 9.780327\times1.0053024 = \qty{9.832186}{m.s^{-2}}$.
(b) $9.832186 - 9.780327 = \qty{0.051859}{m.s^{-2}}$, that is
\qty{5186}{\milli\gal}. (c) $3392/5186 = 0.65$: the centrifugal effect explains
about two thirds of the difference; the flattening, which puts the poles nearer
the centre, the rest.
\end{solution}

\begin{solution}{exo:b1:isostasy-gravity:3}
$r = h\rho_c/(\rho_m - \rho_c) = 3\times2800/500 = \qty{16.8}{km}$. The \omterm{def:b1:layered-earth:layers}{crust} is
thicker by $h + r = \qty{19.8}{km}$ (that is $h\rho_m/(\rho_m - \rho_c) =
3\times3300/500$).
\end{solution}

\begin{solution}{exo:b1:isostasy-gravity:4}
(a) $1/0.3086 = \qty{3.24}{m}$. (b) $0.01/0.3086 = \qty{0.032}{m}$: the heights
must be known to about \qty{3}{cm}, which is why a survey records every
station's height with satellite positioning or levelling.
\end{solution}

\begin{solution}{exo:b1:isostasy-gravity:5}
(a) $\rho = 2700\times100/(100 + h)$: \qty{2673}{kg.m^{-3}} for \qty{1}{km},
\qty{2621}{kg.m^{-3}} for \qty{3}{km}, \qty{2571}{kg.m^{-3}} for \qty{5}{km}.
(b) $\rho = (2700\times100 - 1030\times4)/96 = \num{265880}/96 =
\qty{2770}{kg.m^{-3}}$: the column under the sea is denser. (c) Pratt's model
asks for lateral changes of density at a constant depth. Under the ocean floor
such changes exist: young, hot \omterm{def:b1:layered-earth:lithosphere}{lithosphere} near a ridge is lighter than old,
cold \omterm{def:b1:layered-earth:lithosphere}{lithosphere} far from it (\cref{ch:b1:heat-budget}). Under the continents,
on the contrary, the thickness of the \omterm{def:b1:layered-earth:layers}{crust} is seen to change with the relief.
\end{solution}

\begin{solution}{exo:b1:isostasy-gravity:6}
(a) Station 1: $g - \gamma = 9.807751 - 9.808911 = \qty{-0.001160}{m.s^{-2}} =
\qty{-116.0}{\milli\gal}$; \omterm{def:b1:isostasy-gravity:anomalies}{free-air correction} $0.3086\times350 =
\qty{108.0}{\milli\gal}$, so $\Delta g_{\mathrm{FA}} = \qty{-8.0}{\milli\gal}$;
\omterm{def:b1:isostasy-gravity:anomalies}{Bouguer correction} $0.112\times350 = \qty{39.2}{\milli\gal}$, so
$\Delta g_{\mathrm B} = \qty{-47.2}{\milli\gal}$.
Station 2: $g - \gamma = \qty{-535.9}{\milli\gal}$; \omterm{def:b1:isostasy-gravity:anomalies}{free-air correction}
$0.3086\times1850 = \qty{570.9}{\milli\gal}$, so $\Delta g_{\mathrm{FA}} =
\qty{+35.0}{\milli\gal}$; \omterm{def:b1:isostasy-gravity:anomalies}{Bouguer correction} $0.112\times1850 =
\qty{207.2}{\milli\gal}$, so $\Delta g_{\mathrm B} = \qty{-172.2}{\milli\gal}$.
(b) At station 2 the \omterm{def:b1:isostasy-gravity:anomalies}{Bouguer anomaly} is strongly negative: below the mountains
there is a mass deficit, a root of light \omterm{def:b1:layered-earth:layers}{crust}. The \omterm{def:b1:isostasy-gravity:anomalies}{free-air anomaly}, small
compared with the \omterm{def:b1:isostasy-gravity:anomalies}{Bouguer correction}, shows that the deficit nearly balances
the mountains: they are largely compensated (to about $172/207$, or
\qty{83}{\percent}, if the slab model holds for the whole range).
\end{solution}

\begin{solution}{exo:b1:isostasy-gravity:7}
(a) As in the proof of the proposition, a thin layer at the distance $z$
contributes $G\rho\,2\pi z\,h\int_0^a s\,(s^2 + z^2)^{-3/2}\,\mathrm ds$, now
with the integral stopped at $s = a$:
$\int_0^a s\,(s^2 + z^2)^{-3/2}\,\mathrm ds = 1/z - 1/\sqrt{z^2 + a^2}$. Hence
$g = 2\pi G\rho h\,\bigl(1 - z/\sqrt{z^2 + a^2}\bigr)$, which tends to the slab
value as $a \to \infty$. (b) $z/\sqrt{z^2 + a^2} = 0.1$ gives
$\sqrt{z^2 + a^2} = 10z$, $a = \sqrt{99}\,z = \qty{9.95}{km}$ for
$z = \qty{1}{km}$. (c) The rock within about ten times its depth below the
station supplies \qty{90}{\percent} of the attraction: the terrain correction
needs the relief over a radius of several kilometres to tens of kilometres
around each station, not only beneath it.
\end{solution}

\begin{solution}{exo:b1:isostasy-gravity:8}
(a) $r' = d(\rho_c - \rho_w)/(\rho_m - \rho_c) = 4\times1770/500 =
\qty{14.16}{km}$, and the \omterm{def:b1:layered-earth:layers}{crust} is $35 - 4 - 14.16 = \qty{16.84}{km}$ thick.
(b) Airy's model with one crustal density predicts an oceanic \omterm{def:b1:layered-earth:layers}{crust} more than
three times too thick. The oceanic \omterm{def:b1:layered-earth:layers}{crust} is not thinned continental \omterm{def:b1:layered-earth:layers}{crust}: it
is a different, denser rock, and the two levels of the Earth are two kinds of
\omterm{def:b1:layered-earth:layers}{crust}, as the \omterm{def:b1:isostasy-gravity:hypsometry}{hypsometric curve} suggests.
\end{solution}

\begin{solution}{exo:b1:isostasy-gravity:9}
(a) One second of arc is $\pi/(180\times3600) = 1/\num{206265}$ of a radian:
\qty{5.236}{''} $= \qty{2.538e-5}{rad}$ and \qty{15.885}{''} $=
\qty{7.701e-5}{rad}$. (b) $a_h = g\theta$: $9.79\times\num{2.538e-5} =
\qty{2.485e-4}{m.s^{-2}} = \qty{24.9}{\milli\gal}$ observed, and
$9.79\times\num{7.701e-5} = \qty{75.4}{\milli\gal}$ computed from the visible
mountains. (c) $1 - 5.236/15.885 = 0.670$: the compensation cancelled about two
thirds of the attraction of the visible mountains.
\end{solution}

\begin{solution}{exo:b1:isostasy-gravity:10}
(a) Outside it, a sphere attracts like a point mass at its centre, with
$GM/(x^2 + z^2)$ along the line to the centre; the vertical component is that
times the cosine $z/\sqrt{x^2 + z^2}$, which gives $GMz/(x^2 + z^2)^{3/2}$.
(b) $M = \tfrac43\pi(\qty{2000}{m})^3\times(\qty{-200}{kg.m^{-3}}) =
\qty{-6.70e12}{kg}$, and above the centre $GM/z^2 = \num{6.674e-11}\times
(\num{-6.70e12})/(4000)^2 = \qty{-2.80e-5}{m.s^{-2}} = \qty{-2.80}{\milli\gal}$.
(c) $(1 + x^2/z^2)^{-3/2} = 1/2$ gives $x = z\sqrt{2^{2/3} - 1} = 0.766\,z =
\qty{3.07}{km}$. (d) The half-width at half-maximum of the profile gives the
depth: $z = 1.305\,x_{1/2}$. A deeper body gives a broader, flatter anomaly.
\end{solution}

\begin{solution}{exo:b1:isostasy-gravity:11}
(a) $V = \tfrac13\pi R^2 H = \tfrac13\pi(\qty{15e3}{m})^2(\qty{3e3}{m}) =
\qty{7.07e11}{m^3}$, and $\Delta m = (2800 - 1030)V = \qty{1.25e15}{kg}$.
(b) The centre of mass lies \qty{0.75}{km} above the floor, at
$d = 5 - 0.75 = \qty{4.25}{km}$ below the sea surface. There the potential
anomaly is $T = -G\Delta m/d$, so $N = G\Delta m/(\gamma d) =
\num{6.674e-11}\times\num{1.25e15}/(9.81\times4250) = \qty{2.0}{m}$.
(c) Most of the cone's mass lies in its wide lower part, at 5 to
\qty{16}{km} from the point above the summit: the mean of $1/r$ over the mass
is smaller than $1/\qty{4.25}{km}$, and the point mass overestimates the
potential. Summing thin discs gives about \qty{1.1}{m}.
(d) A root of light \omterm{def:b1:layered-earth:layers}{crust} beneath would put a deficit of the same mass deeper
down; the two potentials nearly cancel, and the bump would be much smaller,
but not zero, since the deficit lies farther from the surface than the excess.
\end{solution}

\begin{solution}{exo:b1:isostasy-gravity:12}
(a) The uplift rate is $-\mathrm dw/\mathrm dt = w_0\,e^{-t/\tau}/\tau = w/\tau$.
(b) $w = 818\,e^{-10000/4000} = 818\times0.0821 = \qty{67.1}{m}$, and the rate is
$\qty{67.1}{m}/\num{4000}\,\unit{yr} = \qty{16.8}{mm.yr^{-1}}$, more than one
and a half times the measured value. (c) With $\tau = 3100$ years, $w =
818\,e^{-3.226} = \qty{32.5}{m}$ and the rate is \qty{10.5}{mm.yr^{-1}}; with
$\tau = \num{68000}$ years, $w = 818\,e^{-0.147} = \qty{706}{m}$ and the rate
is \qty{10.4}{mm.yr^{-1}}. The two imply very different remaining depressions.
An uncompensated depression is a mass deficit of \omterm{def:b1:layered-earth:layers}{mantle} rock: over a broad
region it gives a \omterm{def:b1:isostasy-gravity:anomalies}{free-air anomaly} of about $-2\pi G\rho_m w$, that is
$-0.0419\times3.3\times706 \approx \qty{-98}{\milli\gal}$ for the slow case,
against \qty{-4}{\milli\gal} for the fast one. A gravity survey (or the height
of old shorelines) tells them apart. (d) The time constant measures how fast
the \omterm{def:b1:layered-earth:layers}{mantle} flows back under the depression: its viscosity.
\end{solution}

\begin{solution}{pb:b1:isostasy-gravity:1}
\textbf{1.} $r = h\rho_c/(\rho_m - \rho_c) = 4\times2700/600 = \qty{18}{km}$.

\textbf{2.} The \omterm{def:b1:layered-earth:layers}{crust} is $35 + 4 + 18 = \qty{57}{km}$ thick; its base, the \omterm{def:b1:layered-earth:discontinuity}{Moho},
lies $35 + 18 = \qty{53}{km}$ below sea level.

\textbf{3.} Plateau: \qty{57}{km} of \omterm{def:b1:layered-earth:layers}{crust} and $60 - 53 = \qty{7}{km}$ of \omterm{def:b1:layered-earth:layers}{mantle},
$\num{57000}\times2700 + \num{7000}\times3300 = \num{1.539e8} + \num{2.31e7} =
\qty{1.770e8}{kg.m^{-2}}$. Lowland: $\num{35000}\times2700 + \num{25000}\times3300
= \num{9.45e7} + \num{8.25e7} = \qty{1.770e8}{kg.m^{-2}}$. They are equal, as
\cref{prop:b1:isostasy-gravity:equal-mass} requires.

\textbf{4.} $p = g\times\text{mass per unit area} = 9.81\times\num{1.770e8} =
\qty{1.74e9}{Pa}$, about \qty{1.7}{GPa}, under both columns.

\textbf{5.} The root must make up, by its smaller density, for the mass of the
plateau: each metre of root removes $\rho_m - \rho_c = \qty{600}{kg.m^{-3}}$ of
mass, and each metre of plateau adds \qty{2700}{kg.m^{-3}}, so $2700/600 = 4.5$
metres of root are needed. With a \omterm{def:b1:layered-earth:layers}{crust} of \qty{2900}{kg.m^{-3}}, the ratio
would be $2900/400 = 7.25$: the smaller the contrast, the deeper the root.

\textbf{6.} The \omterm{def:b1:layered-earth:discontinuity}{Moho}, the strong increase in seismic speed, would be found
to deepen under the edge of the plateau, from \qty{35}{km} beneath the lowland
to about \qty{53}{km} beneath the plateau.

\textbf{7.} $\rho = 2700\times100/104 = \qty{2596}{kg.m^{-3}}$.

\textbf{8.} $\num{104000}\times2596.2 = \qty{2.700e8}{kg.m^{-2}}$ and
$\num{100000}\times2700 = \qty{2.700e8}{kg.m^{-2}}$.

\textbf{9.} $2700 - 2596 = \qty{104}{kg.m^{-3}}$, or $104/2700 =
\qty{3.85}{\percent}$.

\textbf{10.} $\Delta\rho/\rho = \alpha\,\Delta T$ gives $\Delta T =
0.0385/\num{3e-5} = \qty{1280}{K}$ on average over the whole \qty{104}{km}
column. That would melt much of the column: heat alone cannot make this plateau
light enough, and Pratt's mechanism cannot be the main one here.

\textbf{11.} Pratt's columns have no root: the \omterm{def:b1:layered-earth:discontinuity}{Moho} stays at about \qty{35}{km}
under the plateau, as under the lowland.

\textbf{12.} \qty{52}{km} is close to Airy's \qty{53}{km} and far from Pratt's
\qty{35}{km}: the plateau is compensated mainly by a crustal root.

\textbf{13.} $0.0419\times2.7\times4000 = \qty{452.5}{\milli\gal}$.

\textbf{14.} Uncompensated (\cref{prop:b1:isostasy-gravity:broad}):
$\Delta g_{\mathrm{FA}} \approx \qty{+452.5}{\milli\gal}$,
$\Delta g_{\mathrm B} \approx 0$.

\textbf{15.} Compensated: $\Delta g_{\mathrm{FA}} \approx 0$ and
$\Delta g_{\mathrm B} \approx \qty{-452.5}{\milli\gal}$. The root's missing mass
gives $-0.0419\times0.6\times\num{18000} = \qty{-452.5}{\milli\gal}$, the same,
since $\Delta\rho\,r = \rho_c h$.

\textbf{16.} $\sin^2\ang{30} = 0.25$ and $\sin^2\ang{60} = 0.75$:
$\gamma = 9.780327\times(1 + 0.0013256 - 0.0000044) = \qty{9.793249}{m.s^{-2}}$.
Then $g - \gamma = 9.781085 - 9.793249 = \qty{-1216.4}{\milli\gal}$; the \omterm{def:b1:isostasy-gravity:anomalies}{free-air
correction} is $0.3086\times4000 = \qty{1234.4}{\milli\gal}$, so
$\Delta g_{\mathrm{FA}} = \qty{+18.0}{\milli\gal}$, and
$\Delta g_{\mathrm B} = 18.0 - 452.5 = \qty{-434.5}{\milli\gal}$.

\textbf{17.} $434.5/452.5 = 0.96$: the plateau is about \qty{96}{\percent}
compensated (the small positive \omterm{def:b1:isostasy-gravity:anomalies}{free-air anomaly} is the uncompensated
\qty{4}{\percent}).

\textbf{18.} No. Near the centre of a broad plateau each anomalous layer acts
as a slab, whose attraction does not depend on its depth: only the mass deficit
per unit area matters, and it is the same, $\rho_c h$, in both models. The two
differ only near the edges, where the depth of the compensating masses changes
the shape of the anomaly.

\textbf{19.} Far from the edge on either side the \omterm{def:b1:isostasy-gravity:anomalies}{free-air anomaly} is near zero.
A few kilometres inside the edge, the plateau beneath the station is close and
extends far on one side only: it still pulls with nearly its full slab value.
The root, \qty{35}{km} or more below, is seen from there as a body ending
almost straight below: it pulls with little more than half its effect. The
difference is a positive peak. A few kilometres outside, the plateau pulls
very little, while the root still pulls with nearly half its effect: a negative
trough. The profile is a positive--negative pair across each edge, as in the
figure of the chapter.

\textbf{20.} The \omterm{def:b1:layered-earth:layers}{crust} is now $57 - 1 = \qty{56}{km}$ thick, $\qty{21}{km}$ more
than the lowland's. In equilibrium this excess splits as $h' + r'$ with
$r' = 4.5\,h'$: $h' = 21/5.5 = \qty{3.818}{km}$ and $r' = \qty{17.18}{km}$.

\textbf{21.} The plateau now stands at \qty{3818}{m}.

\textbf{22.} The rock now at the surface was \qty{1}{km} below the old surface,
at \qty{3000}{m}; it is now at \qty{3818}{m}: it has risen by \qty{818}{m},
that is $1000\times\rho_c/\rho_m$.

\textbf{23.} Each kilometre eroded lowers the surface by $1 - 2700/3300 =
0.182$ of a kilometre, so bringing it down by \qty{4}{km} takes
$4/0.182 = \qty{22}{km}$ of erosion, the whole excess thickness $h + r$. In an
old, worn-down mountain belt the rocks at the surface were therefore once some
twenty kilometres deep: the deep roots of the belt are exposed.

\textbf{24.} The surface comes down by $1000 - 818$: for each kilometre of rock
eroded, the net lowering of the surface is
\textbf{\qty[text-series-to-math=true]{182}{m}}.
\end{solution}
